\documentclass[10pt,a4paper]{amsart}
\usepackage[margin=19mm]{geometry}
\usepackage{amsmath,amssymb,mathtools}
\usepackage[scr=rsfs,cal=euler]{mathalfa}
\usepackage{xfrac}
\usepackage[unicode,hidelinks,bookmarks=false]{hyperref}
\newcommand{\ie}{i.e.\ }
\newcommand{\cf}{cf.\ }
\newcommand{\resp}{respectively\ }
\newcommand{\Subsection}{\subsection}
\providecommand{\nobreakdash}{\nobreak\hbox{-}}
\setlength{\emergencystretch}{2em}
\multlinegap0pt
\usepackage{tikz}
\usepackage[shortlabels]{enumitem}
\newtheorem{lemma}{Lemma}
\title{On the asymptotic behavior of the Repulsive \\ Pressureless Euler-Poisson System}
\author{Nicholas Biglin, Joseph Crachiola, Jack Curtis, Thomas Kunz, Omkar Maralappanavar, Adrian Tudorascu}
\date{Source: arXiv:2601.19733v2 (CC BY 4.0)}
\begin{document}
\maketitle

\noindent\emph{Source and licence.} On the asymptotic behavior of the Repulsive \\ Pressureless Euler-Poisson System. Version arXiv:2601.19733v2 (CC BY 4.0), distributed by arXiv under the Creative Commons Attribution 4.0 International licence, http://creativecommons.org/licenses/by/4.0/, as recorded in the arXiv licence metadata for this exact version and shown on the abstract page https://arxiv.org/abs/2601.19733v2. Full licence notice: \texttt{licenses/arxiv-2601.19733-CC-BY-4.0.txt}.\par\smallskip

\noindent\textit{Editorial scope.} Counted unit: the article's lemma perfect-adjacent-particles (source0.tex lines 954--968). The article's complete original proof of it is delivered with it (source0.tex lines 970--1029, one \texttt{proof} environment of 60 lines). No mathematics is written, repaired or re-derived: the statement and the proof are the article's own lines, in the article's own notation, and no retained line is substituted or re-flowed.  No further source-local context is needed: every cross-reference of every retained line resolves inside the two delivered spans.  Nothing was omitted from any retained span, so the package carries no omission record.  No retained line contains a citation, so the wrapper carries no bibliography and no bibliography entry.  No retained line contains a figure environment, an image-inclusion command or drawing code, so no image is delivered and the package declares no asset dependency.  Editorial formatting only, in addition: an AMS \texttt{amsart} preamble that restates the article's own declarations and macros used by the retained lines, namely the environments \texttt{lemma} on separate counters, together with the standard packages those lines need.  The retained environments print in this document as Lemma 1 in order of appearance, whereas the article prints the same items with the numbering of its own full document.  The complete native source of the article as distributed in the arXiv e-print for this exact version is bundled unchanged as \texttt{sources/193477/source0.tex}.
        \begin{lemma}\label{perfect-adjacent-particles}
            Let $\rho$ be a discrete solution, $\rho_t=\sum_i m_i\delta_{y_i(t)}$, with $y_1(t)\leq y_2(t)\leq\ldots\leq y_n(t)$ for all $t\geq0$. Consider two adjacent particles with masses $m_i$, $m_{i+1}$ and trajectories $y_i(t)$, $y_{i+1}(t)$ which collide with each other but do not collide with other particles beforehand. Then the following statements are equivalent:
           
            (i) The particles collide glancingly.

            (ii) The particles initial positions and velocities satisfy
            \begin{equation}\label{veloc-perfect}
                y_i'(0)-y_{i+1}'(0)=\sqrt{\big(y_{i+1}(0)-y_i(0)\big)(m_i+m_{i+1})}.
            \end{equation}

            (iii) The particles collide at time $t_c$, where 
            \begin{align}\label{TC}
                t_c=2\sqrt{\frac{y_{i+1}(0)-y_i(0)}{m_i+m_{i+1}}}.
            \end{align}
        \end{lemma}

        \begin{proof}
            Let $t_0$ be the time of collision, and let $k = \sum_{j < i}{m_j} - \sum_{j > i + 1}{m_j}$. Then, while $t < t_0$, we have 
            \begin{align}
                y_i''(t) &= \frac{k - m_{i+1}}{2}, \ 
                y_i'(t) = y_i'(0) + \frac{k - m_{i+1}}{2} t, \label{vel__reference_i}\\
                y_i(t) &= y_i(0) + y_i'(0)t + \frac{k - m_{i+1}}{4} t^2, \label{pos__reference_i}\\
                y_{i+1}''(t) &= \frac{k + m_{i}}{2},\ 
                y_{i+1}'(t) = y_{i+1}'(0) + \frac{k + m_{i}}{2} t, \label{vel__reference_i+1}\\
                y_{i+1}(t) &= y_{i+1}(0) + y_{i+1}'(0) t + \frac{k + m_{i}}{4} t^2. \label{pos__reference_i+1}
            \end{align}
            Note that \eqref{pos__reference_i} and \eqref{pos__reference_i+1} apply also at $t=t_0$ due to the continuity of position. Now, assume $i$; that the particles collide glancingly. Since the collision is glancing, we have that $y_i(t_0) = y_{i+1}(t_0)$ and $y_i'(t_0) = y_{i+1}'(t_0)$. By the above equalities, $y'_i(t_0) = y'_{i+1}(t_0)$ implies
            \begin{align}
                y_i'(0) + \frac{k - m_{i+1}}{2} t_0 &= y_{i+1}'(0) + \frac{k + m_{i}}{2} t_0,\nonumber \\
                y_i'(0) - y_{i+1}'(0) &= \frac{m_i + m_{i+1}}{2} t_0 \label{velocity_difference_at_0}
            \end{align}
            and $y_i(t_0) = y_{i+1}(t_0)$ implies
            \begin{align}
                y_i(0) + y_i'(0)t_0 + \frac{k - m_{i+1}}{4} t_0^2 &= y_{i+1}(0) + y_{i+1}'(0)t_0 + \frac{k + m_{i}}{4} t_0^2 \nonumber\\
                \big( y'_i(0) - y'_{i+1}(0) \big) t_0 - \frac{m_i + m_{i+1}}{4}t_0^2 &= y_{i+1}(0) - y_i(0). \label{position_difference_at_0}
            \end{align}
            By substituting \eqref{velocity_difference_at_0} into \eqref{position_difference_at_0}, we get
            \begin{align*}
                \frac{m_i + m_{i+1}}{4}t_0^2 &= y_{i+1}(0) - y_i(0) 
                \\
                t_0 &= 2\sqrt{\frac{y_{i+1}(0) - y_i(0)}{m_i + m_{i+1}}} = t_c.
            \end{align*}
            Substituting this value of $t_0$ into \eqref{velocity_difference_at_0}, we get
            \begin{align*}
                y_i'(0) - y_{i+1}'(0) &= \frac{m_i + m_{i+1}}{2} \cdot 2\sqrt{\frac{y_{i+1}(0) - y_i(0)}{m_i + m_{i+1}}} = \sqrt{\big( y_{i+1}(0) - y_i(0) \big) \left(m_i + m_{i+1} \right)}.
            \end{align*}
            Thus, we have proven that \emph{(i)} implies \emph{(ii)} and \emph{(iii)}. Now, assume \eqref{veloc-perfect} holds.
            We want to show that the particles collide glancingly at $t_c$, that is, $y_i(t_c) - y_{i + 1}(t_c) = 0$, $y_i'(t_c) - y_{i + 1}'(t_c) = 0$. From \eqref{pos__reference_i} and \eqref{pos__reference_i+1}, 
            \begin{align*}
                y_i(t_c) - y_{i + 1}(t_c) 
                % =& \left( y_i(0) + y_i'(0)t_c + \frac{k - m_{i+1}}{4} t_c^2 \right) - \left( y_{i+1}(0) + y_{i+1}'(0)t_c + \frac{k + m_{i}}{4} t_c^2 \right) \\
                =& -\big( y_{i+1}(0) - y_i(0) \big) + \big( y'_i(0) - y_{i+1}'(0)\big) t_c - \left( \frac{m_i + m_{i+1}}{4} \right) t_c^2 \\
                =& -\big( y_{i+1}(0) - y_i(0) \big) + \left( \sqrt{\big( y_{i+1}(0) - y_i(0) \big) \left(m_i + m_{i+1} \right)} \right) \left( 2\sqrt{\frac{y_{i+1}(0) - y_i(0)}{m_i + m_{i+1}}} \right)  \\ 
                &\quad\quad - \left( \frac{m_i + m_{i+1}}{4} \right) \left( 4 \frac{y_{i+1}(0) - y_i(0)}{m_i + m_{i+1}} \right) \nonumber \\
                =& -\big( y_{i+1}(0) - y_i(0) \big) + 2 \big( y_{i+1}(0) - y_{i}(0) \big) - \big( y_{i+1}(0) - y_{i}(0) \big)\ = \ 0,
            \end{align*}
            and from \eqref{vel__reference_i} and \eqref{vel__reference_i+1} we have
            \begin{align*}
                y_i'(t_c) - y_{i + 1}'(t_c) 
                % =& \left( y_i'(0) + \frac{k - m_{i+1}}{2} t_c \right) - \left( y_{i+1}'(0) + \frac{k + m_{i}}{2} t_c \right) \\
                =& \left( y_i'(0) - y_{i+1}'(0) \right) - \frac{m_{i+1} + m_i}{4} t_c \\
                =& \sqrt{\big( y_{i+1}(0) - y_i(0) \big) \left(m_i + m_{i+1} \right)} - \frac{m_{i+1} + m_i}{2} 2\sqrt{\frac{y_{i+1}(0) - y_i(0)}{m_i + m_{i+1}}} 
                % \\ =& \sqrt{\big( y_{i+1}(0) - y_i(0) \big) \left(m_i + m_{i+1} \right)} - \sqrt{\big( y_{i+1}(0) - y_i(0) \big) \left(m_i + m_{i+1} \right)} \\
                = \ 0.
            \end{align*}
            Thus \emph{(ii)} implies \emph{(i)}. Finally, we will prove that \emph{(iii)} implies \emph{(ii)}, and therefore \emph{(i)} as well by the above. Given the assumption that $y_i(t_c) = y_{i+1}(t_c)$ and \eqref{pos__reference_i}, \eqref{pos__reference_i+1}, we argue as follows.
            \begin{align*}
                y_i(0)& + y_i'(0)t_c + \frac{k - m_{i+1}}{4} t_c^2 = y_{i+1}(0) + y_{i+1}'(0) t_c + \frac{k + m_{i}}{4} t_c^2 
                \\ 
                % &\Rightarrow \left (y_i'(0) - y_{i+1}'(0) \right) t_c = y_{i+1}(0) - y_i(0) + \left( \frac{k + m_{i}}{4} - \frac{k - m_{i+1}}{4} \right) t_c^2 \\
                &\Rightarrow \left (y_i'(0) - y_{i+1}'(0) \right) t_c = y_{i+1}(0) - y_i(0) + (m_i + m_{i+1}) \frac{y_{i+1}(0)-y_i(0)}{m_i+m_{i+1}} \\
                &\Rightarrow y_i'(0) - y_{i+1}'(0) = \frac{2\big(y_{i+1}(0) - y_i(0)\big)}{2\sqrt{\frac{y_{i+1}(0)-y_i(0)}{m_i+m_{i+1}}}} 
                % \\&\Rightarrow y_i'(0)-y_{i+1}'(0)
                =\sqrt{\big(y_{i+1}(0)-y_i(0)\big)(m_i+m_{i+1})}. \qedhere
            \end{align*}
        \end{proof}

\end{document}
